% Shared module: the statistical layer (split layer, cloud, Dirichlet-process clustering, decision).
\subsection{Gaussian prior on the components and the split layer}\label{sec:bayes-split}
The core returns factors and components for one field. The first statistical question is how well the
components are determined given the factors. On the observation window write
\begin{equation}\label{eq:linmodel}
  u=\sum_{j=1}^NB_ja_j+\varepsilon=\sumop a+\varepsilon,\qquad \varepsilon\sim\Normal(0,\noise^2I),\qquad a\sim\Normal(0,\omega^2I),
\end{equation}
with $B_j$ the orthonormal kernel bases of \cref{sec:alg-interference} and $a_j$ the coefficients of component $j$ in
that basis. As in \cref{sec:alg-accept}, $\omega^2$ is the prior variance of each coefficient, now in the orthonormal
basis: the prior covariance of $a$ is the diagonal matrix $\omega^2I$. This Gaussian prior on the \emph{coefficients of the components} is the second
prior of the method; the first, in \cref{sec:alg-interp}, is on the field. Given the data, the coefficients are again
Gaussian, centred at the regularized least-squares fit and with covariance $\Gamma_a$:
\begin{equation}\label{eq:postcov}
  a\mid u\ \sim\ \Normal\big(\hat a,\Gamma_a\big),\qquad
  \hat a=\Big(\sumop^H\sumop+\frac{\noise^2}{\omega^2}I\Big)^{-1}\sumop^Hu,\qquad
  \Gamma_a=\noise^2\Big(\sumop^H\sumop+\frac{\noise^2}{\omega^2}I\Big)^{-1}.
\end{equation}
The layer therefore does not return one split of $u$ but a Gaussian distribution of splits around the fitted one.

\begin{proposition}[The angle sets the error bars]\label{prop:bayes-angle}
For two groups the eigenvalues of $\sumop^H\sumop$ are $1\pm\cos\theta_k$, with $\theta_k$ the principal
angles between the sampled kernels, together with the eigenvalue $1$ of multiplicity
$\abs{\dim M_1-\dim M_2}$. In the worst direction orthogonal to $\ker\sumop$ the posterior variance of the component
coefficients is
\begin{equation}\label{eq:worstvar}
  \frac{\noise^2}{\gamma_2+\noise^2/\omega^2},\qquad \gamma_2=1-\cF .
\end{equation}
On the shared directions ($\theta_k=0$, the overlap $\deltaP$) the likelihood is flat, the posterior equals
the prior, and the posterior mean splits the shared part equally.
\end{proposition}

\begin{proof}
With orthonormal $B_1,B_2$, $\sumop^H\sumop=\begin{psmallmatrix}I&B_1^HB_2\\B_2^HB_1&I\end{psmallmatrix}$; the singular
values of $B_1^HB_2$ are the $\cos\theta_k$, which gives the eigenvalues $1\pm\cos\theta_k$ (and $1$ on the
directions of the larger space without a partner), and \eqref{eq:postcov} is diagonal in the same basis. The
smallest nonzero eigenvalue $1-\cF$ gives
\eqref{eq:worstvar}. For $\theta_k=0$ the likelihood depends on $a_1+a_2$ along that direction only, so
the posterior along $a_1-a_2$ is the prior, centred at zero; its variance there is $\omega^2$.
\end{proof}

Equation~\eqref{eq:worstvar} is the Bayesian form of the stability bound~\eqref{eq:stability}: credible
intervals grow like $\noise/\sin(\thetaF/2)$ as the angle closes, and are capped by the prior when
$\thetaF\to0$. The scale $\omega^2$ is not set by hand: it is the value that makes the observed field most probable
(empirical Bayes),
\[
  \hat\omega^2=\operatorname*{arg\,max}_{\omega^2}\ \log\Normal\big(u;\,0,\ \noise^2I+\omega^2\sumop\sumop^H\big),
\]
maximized over a grid from $10^{-6}$ to $10^{2}$. The same evidence compares the model
$u\in M_1\oplus M_2$ with the model with one factor, which gives the probability $P(\text{split})$; as
$\thetaF\to0$ the evidence ratio tends to one and $P(\text{split})$ returns to its prior. The energy shares $w_j$ are nonlinear functions of $a$, so their intervals are
obtained by Monte Carlo, without linearization: the layer draws $400$ coefficient vectors $a^{(1)},\dots,a^{(400)}$ from
the Gaussian~\eqref{eq:postcov}, computes the shares of each, and reports their quantiles.

\subsection{Many samples, one cloud: grouping across samples}\label{sec:bayes-cloud}
The core runs on each of the $\Nsamp$ posterior field samples $u^{(q)}$ of \cref{sec:alg-interp}. Every sample
returns its own factors, with energy shares $w_j^{(q)}$. The same physical group appears once per sample, with slightly different
coefficients, and sometimes split into two near-copies or merged with a neighbour. The second grouping
problem is therefore to recognize the same factor across samples.

Each factor is mapped to a point of a fixed coefficient space: the monic coefficient vector
$\beta\in\C^{m\times(d+1)}$ of its symbol, zero-padded to the largest degrees in play, so that equal
operators give equal vectors whatever their position in a sample's list. The \emph{cloud} is the weighted
point measure
\begin{equation}\label{eq:cloud}
  \cloud=\sum_{q=1}^{\Nsamp}\pi_q\sum_jw_j^{(q)}\,\delta_{\beta_j^{(q)}},
\end{equation}
which gives each factor the mass of its sample times its energy share. Here $\pi_q$ is the weight of the field
sample $u^{(q)}$, inherited from the marginalization over the Gaussian-process hyperparameters~\eqref{eq:marginal-eta};
$w_j^{(q)}$ is the energy share~\eqref{eq:energy} of factor $j$ in that sample, computed from the components of that
sample (one value per factor and sample, not an average over the Monte Carlo draws above); and $\delta_{\beta}$ is a unit
point mass at the coefficient vector $\beta$, so that each factor appears in the cloud as one point. A sharp equation is a cloud with
a few tight clusters; an uncertain one is a diffuse cloud.

\paragraph{Dirichlet-process clustering}
Each group of points in the cloud is modelled as a Gaussian cluster around the coefficient vector of its factor, so
the cloud as a whole is described by a Gaussian mixture model (GMM): a weighted sum of Gaussian densities, one per
group, with weights given by the cluster masses. A finite GMM needs the number of groups in advance. The number of
groups is not known, so we use its Dirichlet-process version, which lets the data choose it: a Dirichlet-process Gaussian mixture
(truncated stick-breaking with eight mixture atoms, diagonal covariances, variational
inference~\cite{BleJor06}), separately for each degree $m$ in $\tau$. The concentration $\alpha$ is chosen
on the grid $\{10^{-3},\dots,10\}$ by the evidence lower bound. Each factor enters the fit once plus a
number of copies proportional to its mass; replicating every factor many times instead splits one group
into several clusters. Two clusters that never occur together in the same sample are merged when a
full-covariance BIC prefers one Gaussian; full covariances are needed because coefficients of one factor
are strongly correlated (in the worked example of \cref{sec:worked}, the constant and the $\xi^2$ coefficient of the diffusion factor
have correlation $0.78$ across 32 samples).

\paragraph{The final distribution}
The output is a distribution, not a point estimate:
\begin{itemize}
  \item the factors (cluster centres) with $90\%$ intervals for every coefficient (quantiles over posterior samples, or over
  parametric-bootstrap refits on gridded input);
  \item the posterior $P(k)$ of the number $k$ of groups, by Monte Carlo over samples and cluster
  responsibilities, together with the posterior over structures from~\eqref{eq:bic};
  \item the energy shares $w_j$ with intervals, $P(\text{split})$, $\gN$, $\cF$ and the distribution of $\rho$;
  \item the label of \cref{sec:alg-labels}.
\end{itemize}
The effective number of groups $\keff=\exp(-\sum_j\bar w_j\log\bar w_j)$, with $\bar w_j$ the cluster masses, is reported
but never rounded to decide $k$: rounding is fragile (a spurious mass of $1.4\%$ moves $\keff$ from $2.00$ to
$2.13$), so $k$ is read from $P(k)$.

\paragraph{When is the model good enough?}
The decision combines four quantities, in this order: the raw-data goodness of fit (otherwise \MIS); the
concentration of $P(\Mfit)$ on one structure (at least $0.9$) and of $P(k)$ on one value; the counting
condition of \cref{prop:counting} (otherwise \IDK); and the conditioning, read through $\gN$ and
$P(\text{split})$. A small Friedrichs angle does not reject the equation, it widens the intervals of the
components and lowers $P(\text{split})$. The reader therefore always sees how well each group is
separated from the others, not only whether the equation fits.
